% Drawing 34/09-1 -- folder 34, batch 09, archive page 167 (PDF page 168).
% Transcription: En marge gauche, un schéma vertical : $\bar X \supset \bar X|U$, au-dessus de $S \supset U$, avec $\xi$ et une flèche vers $\bar X|U$ ; à côté, une note entourée, en partie raturée, où se lisent « remplacer … point … », « dans $S$ », « hypothèse que $I = \pi_1(\ldots)$ opère … ».
% Kind: vertical schema with arrow and circled note. Method 2 (schematic redraw from the facsimile). Status: partial.
% Notes: Bottom of the left margin: X-bar contains X-bar|U above S contains U, joined by vertical strokes; an arrow from the right points to X-bar|U. A long thin pencil arrow rises from the circled note up the margin. The circled note is drawn as its outline only; its words are not reproduced. Where the transcription places xi could not be fixed; it is put at the tail of the arrow.
% Coordinates: 1 unit = 1 mm = 1/100 of the crop width (bbox in page fractions: [0.06, 0.2, 0.3, 0.72]).
% Redrawn by Claude (claude-opus-5-5), issue #73, 2026-11-07; not checked by a human; not compiled (no TeX in the session).
% Rights: reproduces part of the work -- behind RENDER_TRANSCRIPTS, never published by itself.
\documentclass[tikz,border=3pt]{standalone}
\usepackage{amsmath,amssymb}
\usetikzlibrary{patterns,arrows.meta}
\tikzset{pencil/.style={black!45}, struck/.style={black!35}, mark/.style={circle,draw,inner sep=1.3pt}, pt/.style={circle,fill,inner sep=0.9pt}, opt/.style={circle,draw,fill=white,inner sep=1.1pt}}
\begin{document}
\begin{tikzpicture}[x=1mm,y=1mm,line cap=round,line join=round,font=\scriptsize]
\node[inner sep=1pt] at (21,-250) {\scriptsize $\bar X$};
\node[inner sep=1pt] at (30,-250) {\tiny $\supset$};
\node[inner sep=1pt] at (45,-250) {\scriptsize $\bar X|U$};
\draw (22,-258) -- (22,-277);
\draw (40,-258) -- (40,-277);
\node[inner sep=1pt] at (21,-286) {\scriptsize $S$};
\node[inner sep=1pt] at (29,-286) {\tiny $\supset$};
\node[inner sep=1pt] at (38,-286) {\scriptsize $U$};
\draw[-{Stealth[length=2mm]}] (78,-249) -- (57,-249);
\node[inner sep=1pt] at (85,-247) {\tiny $\xi$};
\draw[-{Stealth[length=2mm]},pencil] plot[smooth, tension=0.55] coordinates {(40,-235) (38,-150) (42,-60) (48,-15)};
\draw[pencil] plot[smooth, tension=0.55] coordinates {(8,-150) (40,-140) (75,-160) (80,-200) (60,-235) (20,-235) (8,-200) (8,-150)};
\end{tikzpicture}
\end{document}
